\partdivider{D}{Systems and Rotation}{Real bodies have size, and they spin. This part carries every idea so far over to bodies of finite size, finding the one point that moves simply, and building the rotational twin of every translational law.}{\chref{ch:com} Centre of Mass and System Momentum\\[1.2mm]
\chref{ch:rotation} Torque, Angular Momentum and Moment of Inertia\\[1.2mm]
\chref{ch:rolling} Rolling Motion}

\ncertchapter{13}{Centre of Mass and System Momentum}\label{ch:com}

Every chapter so far has treated bodies as particles: points with mass and no size. A block on an incline,
a car on a banked road, a planet in orbit, each was drawn as a dot. That is a remarkable simplification,
because nothing real is a particle. A thrown spanner tumbles end over end, a diver twists in mid-air, an
exploding shell scatters fragments in every direction. Why should the particle picture have worked at all?

This chapter supplies the answer, and in doing so justifies every dot drawn so far. For any body or system of
bodies there is one point, the centre of mass (CM), that moves exactly as a single particle would if all the mass
were concentrated there and all the external forces acted there. Internal forces, however complicated, cannot
change its velocity. The tumbling spanner's CM follows a clean parabola; the exploding shell's fragments fly
apart, but their CM continues along the original trajectory as though nothing had happened, at least until
the first fragment strikes the ground.

The chapter first locates the CM, for collections of particles and for continuous bodies (a body with a hole
is handled by treating the hole as negative mass). It then proves that the CM moves like a particle, and
finds that the momentum of a system is its total mass times the velocity of its CM.

\section{Locating the Centre of Mass}

We want one point that stands for the whole system. A plain average of positions will not do, because a
heavy part should count for more than a light one, so the centre of mass is a weighted average position, with
each part of the system weighted by its mass.

\begin{defbox}{The centre of mass}
For a collection of particles, labelled $i=1,2,\dots$, with masses $m_i$ (scalars, in \SI{}{kg}) at
position vectors $\vr_i$ (in \SI{}{m}), the centre of mass is the mass-weighted average position:
\[
\vr_{\text{cm}}=\frac{\sum_i m_i\vr_i}{\sum_i m_i}=\frac{1}{M}\sum_i m_i\vr_i,
\qquad\text{or for a continuous body}\quad \vr_{\text{cm}}=\frac{1}{M}\int \vr\,dm,
\]
where $M=\sum_i m_i$ is the total mass. (For a continuous body the sum becomes an integral over tiny
pieces of mass $dm$, with $M=\int dm$; the interlude below shows how to set one up.) $\vr_{\text{cm}}$ is a position vector, so each
coordinate is handled separately: $x_{\text{cm}}=\frac{1}{M}\sum m_i x_i$, and likewise for $y,z$.
\end{defbox}

For masses $m_1$ and $m_2$ a distance $d$ apart, with the origin at $m_1$, it gives
$x_{\text{cm}}=m_2d/(m_1+m_2)$: the midpoint for equal masses, and otherwise nearer the heavier one, at the point where the pair would balance (Figure~\ref{fig:cmbalance}). Moving
the origin shifts every $\vr_i$, and $\vr_{\text{cm}}$, by the same vector, so the point itself does not
depend on where we put the axes.

\begin{figure}[htbp]
\centering
\tikzfig{cm_balance}
\figcaption{fig:cmbalance}{The centre of mass is the balance point. Its distances to the two bodies are in inverse proportion to their masses, so three times the mass means one-third of the distance, exactly like the pivot of a balanced seesaw.}
\end{figure}

\subsection{Where the centre of mass actually is}

The definition tells us how to compute the point; before computing, it pays to know where to expect it.
Symmetry does most of the work: for a uniform body with a centre of symmetry the CM sits at that geometric centre, and if the body
has a line or plane of symmetry the CM must lie on it (the average of a symmetric distribution cannot prefer one side).
Two facts routinely surprise people:

\begin{trapbox}{two things that catch people out}
\textbf{The CM need not be inside the body.} A ring's CM is at its centre, which is air. A horseshoe's
is outside the metal. Nothing sits at the centre of mass; it is a computed average, not a place where
stuff lives. \\[3pt]
\textbf{The CM is not the centre of gravity (CG) in general.} The CM is where mass averages; the CG is where
gravity effectively acts. They coincide in a \emph{uniform} field, which covers every problem in this book, so
we use them interchangeably. For a body large enough that gravity varies across it (the Moon, say) they
separate slightly, and the resulting torque is what keeps the Moon showing us one face.
\end{trapbox}

\begin{tablebox}{Standard centres of mass}
\centering
\begin{tabular}{@{}ll@{\hspace{2.5em}}ll@{}}
\toprule
Uniform rod & midpoint & Triangular lamina & centroid, $h/3$ above the base\\
Disc, ring, sphere & centre & Solid cone & $h/4$ from the base\\
Semicircular ring & $2R/\pi\approx0.637R$ & Solid hemisphere & $3R/8=0.375R$ from flat face\\
Semicircular disc & $4R/3\pi\approx0.424R$ & Hemispherical shell & $R/2$ from flat face\\
\bottomrule
\end{tabular}
\end{tablebox}

For the semicircles and hemispheres the distance is measured from the centre of the full circle or sphere,
along the axis of symmetry. Each entry is derived below.

\subsection{Interlude: how to set up a mass integral}

The table lists results; to \emph{derive} them we need the integral form of the definition, and every such
derivation is the same three moves. The same moves will also give the \emph{moment of inertia} of
\chref{ch:rotation} (a measure of how hard a body is to spin, built from $r^2\,dm$ with $r$ the distance
from an axis), so this interlude collects them once for both.

\begin{methodbox}{The recipe for any continuous CM or moment-of-inertia integral}
\textbf{1. Slice} the body into infinitesimal pieces $dm$, each small enough that every part of it is at
essentially the same position (for a CM) or the same distance from the axis (for a moment of inertia).
\textbf{2. Express $dm$} through the density and a geometric measure of the slice.
\textbf{3. Integrate} the quantity of interest ($x\,dm$ for a CM coordinate, $r^2\,dm$ for a moment of
inertia) over the whole body, and divide by $M$ for a CM.
\end{methodbox}

\textbf{Densities and the mass element.} Everything hinges on writing $dm$ (a scalar, in kilograms) in
terms of a coordinate we can integrate over, and which density you use depends on the body's dimension:
\begin{itemize}
\item \textbf{A wire or rod} (1-D): linear density $\lambda=\dfrac{dm}{d\ell}$ (units \SI{}{kg/m}), where
$d\ell$ is the length of the piece, so $dm=\lambda\,d\ell$. Uniform $\Rightarrow\lambda=M/\ell$, with
$\ell$ the total length.
\item \textbf{A sheet or lamina} (2-D): surface density $\sigma=\dfrac{dm}{dA}$ (units \SI{}{kg/m^2}), so
$dm=\sigma\,dA$ for a piece of area $dA$. Uniform $\Rightarrow\sigma=M/A$, with $A$ the total area.
\item \textbf{A solid} (3-D): volume density $\rho=\dfrac{dm}{dV}$ (units \SI{}{kg/m^3}), so
$dm=\rho\,dV$ for a piece of volume $dV$. Uniform $\Rightarrow\rho=M/V$, with $V$ the total volume.
\end{itemize}
A density need not be constant: a rod made thicker towards one end has a $\lambda$ that varies along it,
and then $dm=\lambda(x)\,dx$ with $\lambda$ kept inside the integral (Example~\ref{ex:nonuniformrod}).

\textbf{Choosing the slice.} The whole art is picking a slice shape that makes the integrand constant
across it, so the calculus stays one-dimensional: short lengths $dx$ for a rod ($dm=\lambda\,dx$); thin
\emph{rings} of radius $r$ and width $dr$ for a disc spun about its axis (area $dA=2\pi r\,dr$,
circumference times width, every point at the same distance $r$ from the axis); thin \emph{discs} of
thickness $dz$ stacked along the axis for a solid sphere or cone; and rings set by an angle from the axis
for a spherical shell. The derivations below show each in action.

\textbf{Symmetry first.} The CM lies on every axis of symmetry (each mass element has a mirror twin whose
contribution cancels), so you only ever integrate along a direction the symmetry does \emph{not} fix: for a
semicircle, just the axis perpendicular to its straight edge.

\begin{tablebox}{The handful of standard integrals used in this book}
\centering
\[
\int_0^{a}x^n\,dx=\frac{a^{n+1}}{n+1},\qquad
\int_0^{\pi}\sin\theta\,d\theta=2,\qquad
\int_0^{\pi}\sin^3\theta\,d\theta=\frac{4}{3}.
\]
Every CM and moment of inertia in this volume reduces to one of these.
\end{tablebox}

\begin{exbox}{A rod whose density grows along it}\label{ex:nonuniformrod}
A rod of length $\ell$ lies along the $x$-axis from $0$ to $\ell$, and its linear density rises steadily from
zero: $\lambda=kx$, with $k$ a constant (units \SI{}{kg/m^2}). Where is its CM? Symmetry is no help, since one
end is heavier, so we follow the recipe: slice into lengths $dx$, write $dm=\lambda\,dx=kx\,dx$, and integrate:
\[
M=\int_0^\ell kx\,dx=\frac{k\ell^2}{2},\qquad
x_{\text{cm}}=\frac{1}{M}\int_0^\ell x\,(kx\,dx)=\frac{k\ell^3/3}{k\ell^2/2}=\frac{2\ell}{3}.
\]
The CM lies beyond the midpoint, towards the dense end, as it must; and $k$ cancels, so only \emph{how}
the density varies matters, not its overall size. \ans{$x_{\text{cm}}=2\ell/3$}
\end{exbox}

\subsection{Deriving the standard centres of mass}

With the recipe tested on the rod, we can now derive each entry in the table. In every case symmetry fixes
all but one coordinate, so only a single integral remains. We measure each from the natural origin and quote the
result as a fraction of the size.

\begin{derivbox}{Triangular lamina (centroid at $h/3$)}
Take a triangle of height $h$ with its base at $y=0$ and apex at $y=h$. Slice it into thin strips
parallel to the base at height $y$, thickness $dy$ (Figure~\ref{fig:cmdtri}). By similar triangles the strip's width shrinks
linearly from its base value $b$ to zero at the apex: $\text{width}=b\left(1-\dfrac{y}{h}\right)$. With
uniform surface density $\sigma$, the strip's mass is $dm=\sigma\,b\left(1-\dfrac{y}{h}\right)dy$, and
\[
y_{\text{cm}}=\frac{\int y\,dm}{\int dm}
=\frac{\displaystyle\int_0^h y\left(1-\tfrac{y}{h}\right)dy}{\displaystyle\int_0^h\left(1-\tfrac{y}{h}\right)dy}
=\frac{\tfrac{h^2}{2}-\tfrac{h^2}{3}}{h-\tfrac{h}{2}}=\frac{h^2/6}{h/2}=\frac{h}{3}.
\]
The centroid sits at height $h/3$ above the base. The same argument holds with any side taken as the
base, so the centroid lies one-third of the way along each median, measured from the midpoint of the
side.
\end{derivbox}

\Needspace*{23\baselineskip}
\begin{figure}[htbp]
\centering
\tikzfig{cmd_triangle}
\figcaption{fig:cmdtri}{The slice for the triangle: a strip of thickness $dy$ at height $y$, whose width
shrinks linearly to zero at the apex. Each strip's midpoint lies on the median from the apex, so only its
height $y$ enters the average.}
\end{figure}

\begin{derivbox}{Semicircular wire (ring), CM at $2R/\pi$}
A uniform semicircular \emph{wire} of radius $R$ sits with its diameter along the $x$-axis, curving up
into $y>0$. By symmetry $x_{\text{cm}}=0$; we need $y_{\text{cm}}$. Parametrise a point on the wire by the
angle $\theta$ from the $x$-axis (Figure~\ref{fig:cmdwire}): its height is $y=R\sin\theta$, and a small arc subtends
$d\ell=R\,d\theta$, so with linear density $\lambda=M/(\pi R)$ the element is $dm=\lambda R\,d\theta$
(the same for every $\theta$). Then
\[
y_{\text{cm}}=\frac{\int_0^\pi (R\sin\theta)\,\lambda R\,d\theta}{\int_0^\pi \lambda R\,d\theta}
=\frac{R\int_0^\pi \sin\theta\,d\theta}{\int_0^\pi d\theta}
=\frac{R\cdot 2}{\pi}=\frac{2R}{\pi}\approx0.637R.
\]
The single standard integral $\int_0^\pi\sin\theta\,d\theta=2$ does all the work.
\end{derivbox}

\begin{figure}[htbp]
\centering
\tikzfig{cmd_semiwire}
\figcaption{fig:cmdwire}{The slice for the semicircular wire: a short arc $d\ell=R\,d\theta$ at angle
$\theta$, sitting at height $y=R\sin\theta$. Sweeping $\theta$ from $0$ to $\pi$ covers the whole wire,
and every element has the same mass $\lambda R\,d\theta$.}
\end{figure}

\begin{derivbox}{Semicircular lamina (disc), CM at $4R/3\pi$}
Now fill the semicircle in. Straight strips parallel to the diameter would work but lead to an awkward
integral, so slice instead into concentric semicircular
\emph{arcs} of radius $r$ (from $0$ to $R$) and thickness $dr$ (Figure~\ref{fig:cmddisc}). Each such arc is itself a semicircular
wire, whose own CM sits at height $\dfrac{2r}{\pi}$ (just derived) and whose mass is
$dm=\sigma\,\pi r\,dr$ (length $\pi r$ times width $dr$). Then
\[
y_{\text{cm}}=\frac{\int_0^R \frac{2r}{\pi}\,(\sigma\pi r\,dr)}{\int_0^R \sigma\pi r\,dr}
=\frac{2\int_0^R r^2\,dr}{\pi\int_0^R r\,dr}
=\frac{2\cdot R^3/3}{\pi\cdot R^2/2}=\frac{4R}{3\pi}\approx0.424R.
\]
Building the disc out of arcs we already understand is quicker than a fresh double integral.
\end{derivbox}

\begin{figure}[htbp]
\centering
\tikzfig{cmd_semidisc}
\figcaption{fig:cmddisc}{The slice for the semicircular disc: nested semicircular arcs. Each arc of radius
$r$ is itself the wire we just solved, with its own centre of mass at $2r/\pi$, so we integrate those
known heights, weighted by each arc's mass, rather than starting over.}
\end{figure}

\begin{derivbox}{Solid cone, CM at $h/4$ from the base}
Stand a solid cone of height $h$ and base radius $R$ on its base and measure $z$ \emph{down from the apex}. At depth $z$ the
cone's radius has grown linearly to $\dfrac{R}{h}z$ (Figure~\ref{fig:cmdcone}), so the thin disc there has volume
$dV=\pi\left(\dfrac{R}{h}z\right)^2 dz\propto z^2\,dz$, giving $dm\propto z^2\,dz$. Then
\[
z_{\text{cm}}=\frac{\int_0^h z\cdot z^2\,dz}{\int_0^h z^2\,dz}=\frac{h^4/4}{h^3/3}=\frac{3h}{4}
\quad\text{from the apex,}
\]
which is $h-\tfrac{3h}{4}=\dfrac{h}{4}$ from the base. The mass piles up towards the wide base, so the CM
sits low, three-quarters of the way down from the tip.
\end{derivbox}

\begin{figure}[htbp]
\centering
\tikzfig{cmd_cone}
\figcaption{fig:cmdcone}{The slice for the cone: a thin disc of thickness $dz$ at depth $z$ below the
apex, whose radius grows linearly as $\tfrac{R}{h}z$. Because the disc's area (and so its mass) grows as
$z^2$, the wide discs near the base dominate the average and pull the CM down.}
\end{figure}

\begin{derivbox}{Solid hemisphere, CM at $3R/8$}
Rest a solid hemisphere on its flat face and measure $z$ up from that face along the axis. A thin disc at
height $z$ (Figure~\ref{fig:cmdhemi}, left) has radius $\sqrt{R^2-z^2}$ (from the sphere's equation), so
$dm=\rho\,\pi(R^2-z^2)\,dz$. Then
\[
z_{\text{cm}}=\frac{\int_0^R z\,(R^2-z^2)\,dz}{\int_0^R (R^2-z^2)\,dz}
=\frac{\big[\tfrac{R^2z^2}{2}-\tfrac{z^4}{4}\big]_0^R}{\big[R^2z-\tfrac{z^3}{3}\big]_0^R}
=\frac{R^4/4}{2R^3/3}=\frac{3R}{8}=0.375R.
\]
\end{derivbox}

\begin{figure}[htbp]
\centering
\tikzfig{cmd_hemisphere}
\figcaption{fig:cmdhemi}{Slicing a hemisphere. Left, the solid: a horizontal disc of thickness $dz$ at height $z$, whose radius $\sqrt{R^2-z^2}$ comes from the sphere's equation; the centre of mass is at $3R/8$. Right, the thin shell: a ring at polar angle $\theta$ of slant width $R\,d\theta$; equal heights carry equal area, so the centre of mass rises to $R/2$.}
\end{figure}

\begin{derivbox}{Hemispherical shell, CM at $R/2$}
A thin \emph{shell} (surface only). Slice it into rings by the polar angle $\theta$ from the axis
(Figure~\ref{fig:cmdhemi}, right). A ring
at angle $\theta$ sits at height $z=R\cos\theta$, has radius $R\sin\theta$, and its slant width is
$R\,d\theta$, so its area is $dA=2\pi(R\sin\theta)(R\,d\theta)$ and $dm=\sigma\,dA\propto\sin\theta\,d\theta$.
Since $z=R\cos\theta$, $dz=-R\sin\theta\,d\theta$, so $dA=2\pi R\,|dz|$: equal slices \emph{along the
axis} carry equal surface area (Archimedes' hat-box result). Hence $dm\propto dz$, uniform in $z$, and
\[
z_{\text{cm}}=\frac{\int_0^R z\,dz}{\int_0^R dz}=\frac{R^2/2}{R}=\frac{R}{2}.
\]
The shell's CM sits higher than the solid hemisphere's ($3R/8$), because a shell has no mass bunched near
the flat base to pull the average down.
\end{derivbox}

\subsection{Composite bodies and the negative-mass trick}

The integrals above handle whole bodies. A body with a piece removed is awkward to integrate over directly,
but there is no need: the CM of the complete body and of the missing piece are both known, and the missing
piece can be treated as a body of negative mass superimposed on the complete one.

\begin{derivbox}{The disc with a hole}
A uniform disc of radius $R$ has a circular hole of radius $R/2$ cut out, the hole's edge just touching
the centre (Figure~\ref{fig:cmhole}). Where is the CM of what remains?
\stepn{1} \textbf{The trick: a hole is negative mass.} A complete disc $=$ (disc with hole) $+$ (the
removed plug). Rearranged, \emph{disc with hole $=$ whole disc $-$ plug}, and the CM formula, being a
sum, lets us enter the plug's mass with a \emph{minus} sign. This is arithmetic, not new physics.
\stepn{2} \textbf{Set up coordinates and masses.} Put the origin at the big disc's centre; the hole's
centre is then at $x=-R/2$. The disc is uniform, so mass goes as area: if the full disc has mass $M$, the
plug (radius $R/2$, hence one-quarter the area) has mass $M/4$, and the remaining piece has mass $3M/4$.
\stepn{3} \textbf{Apply the formula with the plug negative:}
\[
x_{\text{cm}}=\frac{M(0)-\tfrac{M}{4}\!\left(-\tfrac{R}{2}\right)}{M-\tfrac{M}{4}}
=\frac{\tfrac{MR}{8}}{\tfrac{3M}{4}}=\frac{MR}{8}\cdot\frac{4}{3M}=\frac{R}{6}.
\]
\stepn{4} \textbf{Sanity-check.} The CM shifts to $+R/6$, \emph{away} from the hole, which is exactly
right: remove mass from the left and the average moves right. A vanishing hole would give $x_{\text{cm}}\to0$.
\end{derivbox}

\begin{figure}[htbp]
\centering
\tikzfig{cm_hole}
\figcaption{fig:cmhole}{The holed disc is the whole disc minus the plug, so its centre of mass shifts away from the missing piece, to $+R/6$. The same idea handles an L-bracket, a plate with bolt-holes, or any body with a piece removed.}
\end{figure}

\section{The Motion of the Centre of Mass}

We can now locate the CM of almost anything, but the chapter's opening question is still open: why should
that point move like a particle? The theorem proved here answers it, whatever the parts of the system are doing.

\begin{derivbox}{The centre-of-mass theorem}
Start from the definition $M\vr_{\text{cm}}=\sum_i m_i\vr_i$ and differentiate twice with respect to
time, holding every $m_i$ fixed (see the caution on variable mass below). With $\vv_i=d\vr_i/dt$ and
$\va_i=d\vv_i/dt$ the velocity and acceleration of particle $i$, and $\vv_{\text{cm}}$, $\va_{\text{cm}}$
those of the CM,
\[
M\vv_{\text{cm}}=\sum_i m_i\vv_i,\qquad M\va_{\text{cm}}=\sum_i m_i\va_i=\sum_i \vF_i.
\]
The last step is Newton's second law for each particle (in an inertial frame), with $\vF_i$ the net force
on particle $i$. Split every $\vF_i$ into \emph{external} forces (from bodies outside the system) and
\emph{internal} ones (from other particles of the system). The internal forces come in third-law pairs: if
particle $j$ pushes $i$ with $\vF_{ij}$ (first subscript: the particle acted on), then $i$ pushes $j$ with
$\vF_{ji}=-\vF_{ij}$. Summed over the whole system they
cancel pair by pair, leaving only the total external force $\vF_{\text{ext}}$:
\[
\boxed{M\va_{\text{cm}}=\vF_{\text{ext}}.}
\]
\end{derivbox}

\begin{keybox}{The CM moves as a single particle}
The centre of mass accelerates as though \emph{all} the mass were concentrated there and \emph{all} the
external force applied there. Internal forces, however complicated, cannot accelerate the CM. Nor does it
matter \emph{where} on the body an external force acts: a push at the end of a floating plank gives its CM
the same acceleration as the same push at its middle. The point of application decides only how the body
turns about its CM, which is the business of \chref{ch:rotation}.
\end{keybox}

\begin{trapbox}{internal forces cannot move a system}
If the engine is part of the car, how does the car accelerate? The tyres push backwards on the road, and
the road's friction pushes forwards on the tyres: that \emph{external} force moves the CM (\chref{ch:rolling}).
On frictionless ice the same engine only spins the wheels.\\[3pt]
\textbf{Variable mass.} The theorem needs a fixed set of particles. It fails for a rocket alone, whose mass
falls as it burns fuel, but holds for the rocket \emph{together with} its expelled exhaust, whose CM in deep
space stays at rest while the rocket races ahead.
\end{trapbox}

\begin{figure}[htbp]
\centering
\tikzfig{cm_hammer}
\figcaption{fig:cmhammer}{A thrown hammer tumbles end over end, and no point traces a simple path
except the centre of mass, which (air resistance neglected) follows \chref{ch:projectiles}'s parabola,
because gravity is the only external force. This is the retroactive licence for every ``block of mass $m$''
drawn as a dot in the earlier chapters.}
\end{figure}

A thrown hammer shows the theorem at its plainest: however the head and handle tumble, the CM alone
traces a clean parabola (Figure~\ref{fig:cmhammer}). The exploding shell of the chapter's opening is the same theorem, and it shows how much work the theorem
saves: the fragments fly everywhere, but tracking the CM alone locates them.

\begin{exbox}{A shell that explodes at the top of its path}\label{ex:explode13}
A shell fired over level ground would land a distance $D$ from the gun. At the top of its path it explodes
into two equal fragments, one of which is left momentarily at rest and drops vertically. Where does the
other land?
\stepn{1} \textbf{The CM ignores the explosion.} Its forces are internal, so the CM follows the original
parabola to $D$, provided no fragment lands first (the ground's push would be an external force).
\stepn{2} \textbf{Check the landing times.} At the top the shell, of mass $M$, moves horizontally at speed
$u$. Gravity is vertical, so horizontal momentum is conserved through the explosion (\chref{ch:collisions}):
$Mu=\tfrac{M}{2}(0)+\tfrac{M}{2}v_2$ gives $v_2=2u$, also horizontal. Both fragments start
at the same height with no vertical velocity, so they land together and Step~1 applies.
\stepn{3} \textbf{Locate the second fragment.} The first lands directly below the top, at $D/2$. The CM
lands at $D$: $\tfrac{M}{2}\cdot\tfrac{D}{2}+\tfrac{M}{2}x_2=MD$, so $x_2=\tfrac{3D}{2}$.
\ans{$3D/2$ from the gun}
\end{exbox}

\section{The Momentum of a System}

In the explosion we used the CM theorem and momentum conservation side by side, and they agreed: the
fragment's $v_2=2u$ carries it twice as far as the shell would have travelled from the top, from $D/2$ to $3D/2$. That is no accident. The centre-of-mass theorem has an immediate corollary for the total
momentum of a system, and with it a second route to momentum conservation.

\begin{keybox}{System momentum}
The total momentum of a system is that of a single particle of mass $M$ moving with the CM:
$\vc{P}=\sum_i m_i\vv_i=M\vv_{\text{cm}}$. Combined with the CM theorem, $\dfrac{d\vc{P}}{dt}=\vF_{\text{ext}}$,
so \textbf{if the net external force is zero, the total momentum (and hence $\vv_{\text{cm}}$) is
constant.} Conservation of momentum and the CM theorem are one result with two names. \chref{ch:collisions}
proved conservation for two colliding bodies; here it holds for any number. If only one component of
$\vF_{\text{ext}}$ vanishes (no horizontal force, say), only that component of $\vc{P}$ is conserved.
\end{keybox}

The keybox's last line is often the one that matters: a person walking on a boat on still water feels no
horizontal external force, so the CM of person and boat cannot move horizontally, and the boat must slide
back beneath the walker (Figure~\ref{fig:cmboat}).

\begin{exbox}{Walking on a boat}
A \SI{60}{kg} person stands at one end of a \SI{120}{kg} boat of length \SI{3}{m} on frictionless water,
then walks to the other end. How far does the boat move?
\stepn{1} \textbf{Spot the conserved quantity.} No external horizontal force acts, and the system starts
at rest, so horizontally $\vc{P}=0$ throughout: $\vv_{\text{cm}}=0$, and \emph{the CM never moves.}
\stepn{2} \textbf{Write the constraint.} Let the boat slide back a distance $s_{\text{b}}$ while the person advances
$s_{\text{p}}$ relative to the water (both taken as positive distances). Relative to the boat they walked its full length, so
$s_{\text{p}}+s_{\text{b}}=\SI{3}{m}$.
\stepn{3} \textbf{Use the fixed CM.} Since the CM does not move, the mass-weighted displacements must
cancel; the two move in opposite directions, so $m_{\text{person}}\,s_{\text{p}}=m_{\text{boat}}\,s_{\text{b}}$, that is, $60\,s_{\text{p}}=120\,s_{\text{b}}$,
so $s_{\text{p}}=2s_{\text{b}}$.
\stepn{4} \textbf{Solve.} Substituting into $s_{\text{p}}+s_{\text{b}}=\SI{3}{m}$ gives $3s_{\text{b}}=\SI{3}{m}$, so
$s_{\text{b}}=\SI{1}{m}$ and $s_{\text{p}}=\SI{2}{m}$. The
boat slides back \SI{1}{m}, and the person advances only \SI{2}{m}, not the \SI{3}{m} they walked: the
missing metre is the boat retreating beneath them. \ans{boat \SI{1}{m} back, person \SI{2}{m} forward}
\end{exbox}

\begin{figure}[htbp]
\centering
\tikzfig{cm_boat}
\figcaption{fig:cmboat}{With no external horizontal force the centre of mass cannot move, so the boat slides back as the person walks forward, by distances in inverse proportion to the masses (drawn to scale for the worked example). A rocket, with nothing external to push against, works the same way: it throws part of itself backwards.}
\end{figure}

\subsection{The centre-of-mass frame}\label{sec:cmframe}

With no external force the CM moves at constant velocity, so a frame travelling with it is inertial, and
in it the total momentum is zero. In this \textbf{centre-of-mass frame} colliding or exploding pieces always
move with momenta that cancel. It also answers a question raised in \chref{ch:collisions}: how much of a
system's kinetic energy can internal forces change?

\begin{derivbox}{Kinetic energy splits about the CM}
Write each particle's velocity as $\vv_i=\vv_{\text{cm}}+\vv_i'$, where the primed $\vv_i'$ is its
velocity relative to the CM. Then
\[
K=\sum_i\tfrac12m_i(\vv_{\text{cm}}+\vv_i')\cdot(\vv_{\text{cm}}+\vv_i')
=\tfrac12Mv_{\text{cm}}^2+\vv_{\text{cm}}\cdot\sum_i m_i\vv_i'+\sum_i\tfrac12m_iv_i'^2 .
\]
The middle sum is the total momentum seen from the CM frame, $\sum_i m_i\vv_i-M\vv_{\text{cm}}=\vc0$, so
\[
\boxed{K=\tfrac12Mv_{\text{cm}}^2+K'},\qquad K'=\sum_i\tfrac12m_iv_i'^2,
\]
with $K'$ the kinetic energy measured in the CM frame.
\end{derivbox}

For an isolated system momentum fixes the first term, so only $K'$ can be dissipated in a collision or
supplied by an explosion. This
generalises the two-body split of \S\ref{sec:inelastic12}, and it returns in \chref{ch:rolling}, where a
rolling wheel's kinetic energy is translation of its CM plus rotation about it.

\section{Summary}\label{sec:summary13}
The centre of mass is the mass-weighted average position of a system. It moves as a single particle carrying
the total mass and acted on by the total external force, so internal forces cannot change its motion. The
momentum of a system is its total mass times the velocity of its centre of mass, and if the net external force
is zero, both are constant.

\begin{itemize}
\item $\vr_{\text{cm}}=\dfrac1M\sum m_i\vr_i$, or $\dfrac1M\int\vr\,dm$ for a continuous body; each coordinate
      is computed separately.
\item The centre of mass need not lie within the body: a ring's is at its empty centre.
\item Standard results: triangular lamina $h/3$ from the base; semicircular ring $2R/\pi$; semicircular disc
      $4R/3\pi$; solid cone $h/4$ from the base; solid hemisphere $3R/8$; hemispherical shell $R/2$.
\item A body with a hole is a full body plus a negative-mass body occupying the hole.
\item $M\va_{\text{cm}}=\vF_{\text{ext}}$ and $\vc{P}=M\vv_{\text{cm}}$; internal forces (an engine, an
      explosion, a person walking on a boat) cannot move the CM, and the point where an external force acts does not affect $\va_{\text{cm}}$.
\item $K=\tfrac12Mv_{\text{cm}}^2+K'$: only the kinetic energy $K'$ in the CM frame can be changed by internal forces.
\end{itemize}

\section{Problem Types}\label{sec:types13}
Five problem types cover the centre-of-mass questions examiners set. The last two most reward a clear
statement of principle: on the boat, the CM cannot move; in a cut-out body, a hole is negative mass.

\archq{Centre of mass of point masses on a line}
{Several point masses at given positions along a line; find the CM coordinate. A frequent NEET
opening question.}
{Apply $x_{\text{cm}}=\dfrac{\sum m_i x_i}{\sum m_i}$ per axis. Putting the origin at one of the masses removes a term from the numerator.}
{Forgetting to divide by the \emph{total} mass, or mixing up which mass sits where. All positions (and all masses) must be in the same units
before you add.}
{Masses \SI{300}{g}, \SI{500}{g}, \SI{400}{g} sit at $x=0$, \SI{40}{cm}, \SI{70}{cm}. The CM is at:}
{\SI{30}{cm}}{\SI{40}{cm}}{\SI{45}{cm}}{\SI{50}{cm}\archans{B}{$x_{\text{cm}}=\dfrac{300(0)+500(40)+400(70)}{300+500+400}=\dfrac{0+20000+28000}{1200}=\dfrac{48000}{1200}=\SI{40}{cm}$.
Note the \SI{300}{g} at the origin still counts in the denominator even though it adds nothing to the
numerator.}}

\archq{Centre of mass of point masses in a plane}
{``Find the CM of three or more masses at given coordinates in a plane.''}
{Apply the weighted average \emph{per coordinate}:
$x_{\text{cm}}=\dfrac{\sum m_ix_i}{\sum m_i}$ and likewise for $y$. For a symmetric arrangement, put the
origin on the line of symmetry, which fixes one coordinate without calculation.}
{Averaging the positions without weighting by mass. For two bodies, the CM always lies nearer the heavier one.}
{Masses \SI{1}{kg}, \SI{2}{kg} and \SI{3}{kg} sit at $(0,0)$, $(2,0)$ and $(0,2)$ metres. The centre of
mass is at:}
{$(0.67,1.0)$}{$(1.0,0.67)$}{$(0.67,0.67)$}{$(1.0,1.0)$\archans{A}{
$x_{\text{cm}}=\dfrac{1(0)+2(2)+3(0)}{6}=\dfrac46=0.67$;
$y_{\text{cm}}=\dfrac{1(0)+2(0)+3(2)}{6}=\dfrac66=1.0$.}}

\archq{Motion of the CM under internal forces only}
{Two bodies interact (mutual attraction, explosion, a spring, one throwing another) with \emph{no}
external force; asked for the velocity or displacement of the CM, or of one piece.}
{Internal forces cannot move the CM. If the system started at rest, $\vv_{\text{cm}}=0$ \emph{forever}:
the momenta stay equal and opposite, $m_1v_1=m_2v_2$, and for displacements $\sum m_i\Delta x_i=0$. If it
started moving, the CM coasts at constant velocity through the whole event. With only one external force
component zero, apply this along that axis alone.}
{Thinking the CM speeds up or changes direction during an explosion or attraction. It cannot; only an
external force does that.}
{Two particles, initially at rest, move towards each other under their mutual attraction. When their
speeds are $v$ and $2v$, the speed of their centre of mass is:}
{$2v$}{$1.5v$}{$v$}{$0$\archans{D}{The system started at rest with no external force, so its total momentum
is zero at every instant, and $\vv_{\text{cm}}=\vc{P}/M=0$ throughout, whatever the individual
speeds. (The speeds $v,2v$ only tell you the mass ratio, $m_1=2m_2$.)}}

\archq{The CM cannot move: man walking on a boat}
{``A man walks along a boat or plank floating on still water; how far does the boat move?''}
{Treat the water as exerting no horizontal force (no drag); then the CM stays fixed. If the man of mass $m$ walks a distance $\ell$
relative to the boat of mass $m_{\text{b}}$, the boat slides back by
$\boxed{d=\dfrac{m\ell}{m_{\text{b}}+m}}$, and the man's ground displacement is $\ell-d$.}
{Using the man's displacement relative to the \emph{boat} as his displacement relative to the
\emph{ground}. They differ by exactly the boat's own motion.}
{A \SI{50}{kg} man walks \SI{5}{m} along a \SI{200}{kg} boat on still water. The boat moves:}
{\SI{1}{m}}{\SI{1.25}{m}}{\SI{4}{m}}{\SI{5}{m}\archans{A}{$d=\dfrac{m\ell}{m_{\text{b}}+m}=\dfrac{50(5)}{250}
=\SI{1}{m}$, opposite to the man. (B) sets $m\ell=m_{\text{b}}d$, wrongly taking the full \SI{5}{m} as the man's ground
displacement; (C) is the man's ground displacement, $\ell-d$; (D) is the distance walked relative to the boat.}}

\archq{CM of a body with a portion removed}
{``A disc or square has a hole cut in it; where is the new centre of mass?''}
{Treat the removed piece as \textbf{negative mass}:
$x_{\text{cm}}=\dfrac{m_{\text{whole}}x_{\text{whole}}-m_{\text{hole}}x_{\text{hole}}}{m_{\text{whole}}-m_{\text{hole}}}$.
For a uniform lamina the masses are proportional to \emph{areas}, so the density never needs to be known.}
{Forgetting that the removed mass scales with the \emph{square} of the radius for a disc: a hole of
half the radius removes a quarter of the mass, not half.}
{From a uniform disc of radius $R$, a disc of radius $R/2$ is removed with its edge passing through the
centre. The new CM shifts from the centre by:}
{$R/6$ away from the hole}{$R/6$ towards the hole}{$R/4$ away from the hole}{$R/3$ away from the hole\archans{A}{Taking the hole's
mass as $M/4$ with its centre at $x=+R/2$: $x=\dfrac{M(0)-\frac M4\cdot\frac R2}{M-\frac M4}
=\dfrac{-MR/8}{3M/4}=-\dfrac R6$, that is, $R/6$ \emph{away} from the hole.}}
